\documentclass[reprint,amsmath,amssymb,aps,showkeys]{revtex4-2}
\usepackage{graphicx}
\usepackage{booktabs}
\usepackage{hyperref}

\title{Ranked Experiment Program for the Layered ToE (2026)}

\author{C.~M.~Baird}
\affiliation{Baird--ZoraASI Collaboration}

\date{March 2026}

\begin{document}
\maketitle

\section*{Purpose}
This memo turns the current layered Theory of Everything synthesis into a ranked scientific agenda. The goal is not to list every imaginable experiment, but to state which tests carry the most scientific payoff under the present state of the program.

\section*{Ranking principle}
Experiments are ranked by:
\begin{enumerate}
    \item clarity of observable,
    \item vulnerability to confounds,
    \item decisiveness under null,
    \item interpretability by skeptical physicists,
    \item and usefulness for external replication.
\end{enumerate}

\section*{Rank 1: H2 interferometric visibility suppression}
\textbf{Status:} highest-payoff physical signal.\\
\textbf{Why first:} direct observable, one rate parameter, least dependent on protocol-heavy interpretation.\\
\textbf{Pilot stage:} Phase-0 benchtop optical visibility pilot.\\
\textbf{Flagship stage:} upgraded matter-wave / levitated / higher-coherence platform.\\
\textbf{Positive outcome:} motivates immediate escalation and serious anomaly follow-up.\\
\textbf{Null outcome:} still valuable; sets a hard floor on $\Gamma$.\\
\textbf{Messy outcome:} reveals which nuisance channels dominate and what must be fixed before stronger claims.

\section*{Rank 2: H1 branch-selection / non-QRNG stress-test channel}
\textbf{Status:} hardened simulation-primary protocol with adversarial controls.\\
\textbf{Why second:} scientifically useful for stress-testing the estimator and decision logic, but still more vulnerable to objections about protocol dependence if taken as headline evidence.\\
\textbf{Pilot stage:} simulation-primary battery plus any blinded hardware comparator.\\
\textbf{Positive outcome:} only meaningful if the hardware-backed version survives the same frozen analysis chain.\\
\textbf{Null outcome:} constrains one proposed branch-selection route.\\
\textbf{Messy outcome:} often indicates leakage, estimator fragility, or protocol ambiguity.

\section*{Rank 3: H3 multi-channel consistency}
\textbf{Status:} executable posterior/model-comparison prototype.\\
\textbf{Why third:} this is the inference layer that becomes powerful after H2/H1 are stronger, not before.\\
\textbf{Pilot stage:} current implemented H1+H2+fifth-force posterior.\\
\textbf{Flagship stage:} add collider/cosmology tables and outside replication of the whole bundle.\\
\textbf{Positive outcome:} shows the theory survives a joint-weighted constraint layer.\\
\textbf{Null outcome:} overconstraint can efficiently kill broad parameter regions.\\
\textbf{Messy outcome:} usually means missing channels, dominance by one channel, or hidden flexibility.

\section*{Replication stage}
After the first three internal packages, the next stage is not ``invent more theory'' but:
\begin{itemize}
    \item freeze one external contract,
    \item release one outsider-ready manifest,
    \item obtain one independent rerun of the current internal chain,
    \item and then one external physical pilot attempt.
\end{itemize}

\section*{What this ranking means}
The ranking does \textbf{not} mean H1 or H3 are unimportant. It means:
\begin{itemize}
    \item H2 is the sharpest physical blade,
    \item H1 is the hardest stress-test blade,
    \item H3 is the strongest consistency/inference blade.
\end{itemize}

\section*{Bottom line}
If forced to choose where to spend the next scientific unit of energy:
\begin{enumerate}
    \item run H2 pilot,
    \item keep H1 as the estimator/adversarial integrity check,
    \item then let H3 absorb the results into the cross-channel inference layer.
\end{enumerate}

\end{document}
